\documentclass[11pt,a4paper]{article}
\usepackage[utf8]{inputenc}
\usepackage[T1]{fontenc}
\usepackage{lmodern}
\usepackage[spanish,es-nodecimaldot,es-tabla]{babel}
\usepackage[margin=2.2cm,top=2cm,bottom=2.2cm]{geometry}
\usepackage{graphicx}
\usepackage{booktabs}
\usepackage{amsmath,amssymb}
\usepackage{float}
\usepackage{xcolor}
\usepackage{caption}
\usepackage[inline]{enumitem}
\usepackage{titlesec}
\usepackage{microtype}
\usepackage{hyperref}
\usepackage{xurl}
\setlength{\emergencystretch}{3em}

\hypersetup{colorlinks=true, linkcolor=blue!60!black, citecolor=blue!60!black, urlcolor=blue!60!black}
\titleformat{\section}{\large\bfseries\color{blue!60!black}}{\thesection.}{0.5em}{}
\titleformat{\subsection}{\normalsize\bfseries\color{blue!45!black}}{\thesubsection}{0.5em}{}
\titlespacing*{\section}{0pt}{1.0ex plus .3ex}{0.5ex}
\titlespacing*{\subsection}{0pt}{0.8ex plus .2ex}{0.3ex}
\setlist{itemsep=1pt, topsep=2pt, parsep=0pt}
\captionsetup{font=small, labelfont=bf, skip=3pt}
\setlength{\parskip}{2pt}
\setlength{\textfloatsep}{8pt}
\setlength{\floatsep}{6pt}
\graphicspath{{figures/}}
\newcommand{\ci}[2]{{\footnotesize[#1,\,#2]}}

\title{\vspace{-1.2cm}\textbf{Clasificación Automatizada de Exoplanetas Kepler:\\Priorización de Candidatos para Validación Espectroscópica}\\[2pt]
\large Actividad Sumativa — Unidad 4: Inteligencia Artificial}
\author{Maestría en Inteligencia Artificial — Politécnico Grancolombiano}
\textbf{Integrantes del equipo}\\[0.3cm]
Alejandro Quintero Ruiz\\
Angel Masmela\\
Ginna Lorena Delgado\\
Sandra Johana Leguizamon Vasquez\\[0.3cm]
Maestría en Inteligencia Artificial\\
Politécnico Grancolombiano
}
\date{Noviembre de 2026}

\begin{document}
\maketitle
\vspace{-0.6cm}

\begin{abstract}
\noindent La misión Kepler catalogó 9,564 objetos de interés (KOIs): 2,748 confirmados como planetas, 4,839 falsos positivos y 1,977 que siguen como candidatos. Confirmar un candidato exige tiempo escaso de telescopios terrestres. Se presenta un pipeline de aprendizaje supervisado que estima la probabilidad de que un KOI sea un planeta real a partir de 16 parámetros físicos de la señal, del planeta y de la estrella. Se excluyeron las variables del proceso de \textit{vetting} y las barras de error, porque ambas filtran información de la etiqueta. Algoritmo, hiperparámetros y umbral de decisión se eligieron solo con validación cruzada sobre el conjunto de entrenamiento. El modelo final (LightGBM) alcanza en test \textbf{ROC-AUC = 0.979} [0.973, 0.985] y F1 = 0.907 [0.887, 0.924] (IC 95\,\% por \textit{bootstrap}). Con el umbral operativo recupera el 94\,\% de los planetas observando solo el 39\,\% de los KOIs, un ahorro de alrededor del 61\,\% de tiempo de telescopio.
\end{abstract}

\noindent\textbf{Palabras clave:} aprendizaje supervisado, clasificación binaria, exoplanetas, Kepler, \textit{gradient boosting}, \textit{data leakage}.

%=====================================================================
\section{Contexto del problema y justificación}
La misión \textbf{Kepler} de la NASA (2009--2018) buscó exoplanetas con fotometría de tránsito: caídas periódicas del brillo estelar ($\Delta F/F\sim10^{-5}$--$10^{-2}$) cuando un planeta pasa frente a su estrella \cite{borucki2010}. Cada señal periódica se registra como un KOI. Muchas señales no son planetas, sino \textbf{falsos positivos}: binarias eclipsantes, contaminación de una estrella vecina o artefactos instrumentales. Distinguirlos exige seguimiento con espectroscopía de velocidad radial o imagen de alta resolución en telescopios de 8--10\,m, cuyo tiempo es muy limitado. Hoy quedan casi 2,000 KOIs sin resolver, y el problema es aún más urgente en misiones actuales como TESS, que generan miles de candidatos nuevos.

\textbf{¿Por qué aprendizaje automático?}
\begin{enumerate*}[label=(\roman*)]
\item Hay miles de ejemplos con etiqueta experta (\texttt{koi\_disposition}) obtenida tras el seguimiento, lo que permite un enfoque supervisado.
\item Las variables son heterogéneas, no lineales y con valores faltantes, y su relación con la clase es difícil de codificar con reglas manuales.
\item Los costos de error son asimétricos: perder un planeta real (falso negativo) es más grave que gastar una noche en un falso positivo.
\end{enumerate*}
El objetivo no es reemplazar la confirmación observacional, sino \textbf{ordenar la cola de observación} por probabilidad de éxito.

%=====================================================================
\section{Análisis del problema y del dataset}
\textbf{Tipo de problema:} clasificación binaria supervisada, con $y=1$ para \texttt{CONFIRMED} (clase positiva, planeta real) e $y=0$ para \texttt{FALSE POSITIVE}. Es supervisado porque existe una etiqueta experta. La clase \texttt{CANDIDATE} no tiene etiqueta definitiva, así que se excluye del entrenamiento y se usa después para predicción operativa.

\textbf{Datos:} tabla \textit{Kepler Objects of Interest — cumulative} del NASA Exoplanet Archive \cite{nasa}, descargada vía su API TAP pública. Tiene las mismas 50 columnas que el dataset de Kaggle \texttt{nasa/kepler-exoplanet-search-results}.

\begin{table}[H]
\centering\small
\caption{Estructura del dataset}
\setlength{\tabcolsep}{4pt}\begin{tabular}{lr@{\hspace{0.8cm}}lr}
\toprule
KOIs totales & 9,564 & Columnas originales & 50 \\
CONFIRMED ($y{=}1$) & 2,748 (36.2\,\%) & Eliminadas (IDs, etiqueta, leakage, \texttt{\_err}) & 33 \\
FALSE POSITIVE ($y{=}0$) & 4,839 (63.8\,\%) & Predictoras numéricas / categóricas & 16 / 1 \\
CANDIDATE (sin etiqueta) & 1,977 & Train / test (estratificado 80/20) & 6,069 / 1,518 \\
\bottomrule
\end{tabular}
\end{table}

\vspace{-2pt}
Las predictoras describen la \textbf{señal de tránsito} (\texttt{koi\_period}, \texttt{koi\_depth}, \texttt{koi\_duration}, \texttt{koi\_impact}, \texttt{koi\_model\_snr}), el \textbf{planeta inferido} (\texttt{koi\_prad}, \texttt{koi\_teq}, \texttt{koi\_insol}) y la \textbf{estrella} (\texttt{koi\_steff}, \texttt{koi\_slogg}, \texttt{koi\_srad}, \texttt{koi\_kepmag}, coordenadas). La única categórica es \texttt{koi\_tce\_delivname}, la versión del pipeline de detección. Las clases están moderadamente desbalanceadas (1\,:\,1.76).

%=====================================================================
\section{Enfoque metodológico}
\subsection{Preparación de datos y control de \textit{leakage}}
\begin{enumerate}
\item \textbf{Leakage explícito:} se eliminan \texttt{koi\_pdisposition}, \texttt{koi\_score} y las banderas \texttt{koi\_fpflag\_\{nt,ss,co,ec\}}. Son salidas del proceso de vetting de Kepler \cite{thompson2018}, que usa la misma evidencia que define la etiqueta.
\item \textbf{Leakage sutil (barras de error):} se eliminan las 22 columnas \texttt{*\_err1/2}. Cuando un KOI se confirma, sus parámetros, sobre todo los estelares, suelen refinarse con observaciones de seguimiento, y eso reduce sus incertidumbres. La mediana de \texttt{koi\_steff\_err1} es 108\,K en los confirmados frente a 165\,K en los falsos positivos. Las barras de error reflejan, por tanto, el \textit{proceso de confirmación} y no la física de la señal, y un candidato nuevo todavía no las tiene refinadas. Su efecto se cuantifica en la Sección~5.3.
\item También se eliminan los identificadores (\texttt{rowid}, \texttt{kepid}, \texttt{kepoi\_name}, \texttt{kepler\_name}).
\item \textbf{Variables numéricas:} imputación por mediana, robusta a las colas largas de período y profundidad, y \texttt{StandardScaler}, necesario para la regresión logística e inocuo para los árboles. \textbf{Variable categórica:} imputación por moda y \textit{one-hot encoding}. \textbf{Selección de variables:} \texttt{SelectKBest} con información mutua, tratando $k\in\{10,15,\text{todas}\}$ como hiperparámetro.
\end{enumerate}

\textbf{Protocolo de evaluación.} Todo el preprocesamiento va dentro de un \texttt{Pipeline} de scikit-learn \cite{sklearn}, así que se ajusta solo con los datos de entrenamiento de cada \textit{fold}. Las tres decisiones del modelo (algoritmo, hiperparámetros y umbral de decisión) se toman exclusivamente con \texttt{StratifiedKFold(5)} sobre el 80\,\% de entrenamiento. El 20\,\% de test se usa una única vez, para medir. Los intervalos de confianza del 95\,\% se calculan con \textit{bootstrap} de 1,000 remuestreos del test (percentiles). Las diferencias entre modelos se evalúan con \textit{bootstrap} pareado, es decir, ambos modelos sobre los mismos remuestreos.

\subsection{Selección y justificación del algoritmo}
Se compararon cinco modelos con ROC-AUC como métrica, porque es independiente del umbral y robusta al desbalance. A todos se les aplicó corrección de desbalance (\texttt{class\_weight} o \texttt{scale\_pos\_weight}).

\begin{table}[H]
\centering\small
\caption{Comparación de modelos base (ROC-AUC, validación cruzada 5-fold sobre train)}
\begin{tabular}{lcc@{\hspace{1cm}}lcc}
\toprule
Modelo & Media & Desv. & Modelo & Media & Desv. \\
\midrule
Regresión logística & 0.930 & 0.005 & XGBoost & 0.980 & 0.003 \\
Random Forest (300 árboles) & 0.976 & 0.004 & \textbf{LightGBM} & \textbf{0.980} & \textbf{0.003} \\
HistGradientBoosting & 0.979 & 0.004 & & & \\
\bottomrule
\end{tabular}
\end{table}

Los métodos de \textbf{\textit{gradient boosting} de árboles} (XGBoost \cite{xgboost}, LightGBM \cite{lightgbm}) superan claramente al modelo lineal (+0.05 de AUC). Esto confirma que la frontera de decisión es no lineal, y la elección se apoya en las características del problema:

\textit{Fortalezas.}
\begin{enumerate*}[label=(\roman*)]
\item Los árboles son invariantes a transformaciones monótonas, algo útil porque el período va de horas a cientos de días y la profundidad abarca cuatro órdenes de magnitud.
\item Capturan interacciones físicas sin especificarlas, por ejemplo profundidad $\times$ radio estelar $\rightarrow$ radio planetario. Un objeto con $R_p>15\,R_\oplus$ casi siempre es una binaria eclipsante.
\item Tienen regularización explícita (profundidad, número de hojas, \textit{subsampling}).
\item Ofrecen importancias de variables interpretables.
\end{enumerate*}

\textit{Limitaciones.}
\begin{enumerate*}[label=(\roman*)]
\item Sus probabilidades no están calibradas por defecto.
\item Tienen más hiperparámetros y menor interpretabilidad global que un modelo lineal.
\item No extrapolan fuera del rango de entrenamiento, lo que importa al transferir el modelo a otra misión.
\end{enumerate*}

Se optimizaron XGBoost y LightGBM con \texttt{GridSearchCV} (72 combinaciones cada uno: número de árboles $\{200,400\}$, profundidad $\{4,6,8\}$ u hojas $\{15,31,63\}$, \textit{learning rate} $\{0.03,0.1\}$, \textit{subsample} $\{0.8,1.0\}$ y $k$). Se seleccionó \textbf{LightGBM}, con AUC-CV 0.982 (400 árboles, \texttt{num\_leaves=63}, \texttt{learning\_rate=0.03}, \texttt{subsample=0.8}, las 19 variables tras el \textit{one-hot}), frente a 0.981 de XGBoost. Como se muestra en la Sección~5.1, esa diferencia no es estadísticamente significativa.

\subsection{Umbral de decisión operativo}
Para el uso operativo interesa recuperar al menos el 95\,\% de los planetas. El umbral se eligió con predicciones \textit{out-of-fold} del entrenamiento (\texttt{cross\_val\_predict}), en las que cada KOI de train recibe la probabilidad de un modelo que no lo vio. El umbral más alto con recall \textit{out-of-fold} $\geq0.95$ es \textbf{$P\geq0.309$}, y con él se señala el 39.4\,\% de los KOIs de entrenamiento. El test no interviene en esta elección.

%=====================================================================
\section{Desarrollo técnico e implementación}
La solución está en el notebook \texttt{exoplanet\_classification.ipynb} (Python 3.12, pandas, scikit-learn 1.9, XGBoost 3.4, LightGBM 4.7, matplotlib/seaborn). Tiene 15 secciones comentadas que siguen el ciclo completo: carga, limpieza, \textit{split}, comparación de modelos, búsqueda de hiperparámetros, elección de umbral, evaluación con \textit{bootstrap}, interpretabilidad, experimentos de leakage, errores, curvas de aprendizaje y uso operativo. La arquitectura del modelo es:

\begin{center}\small
\texttt{ColumnTransformer}[num: \texttt{SimpleImputer(median)}$\rightarrow$\texttt{StandardScaler} $\mid$ cat: \texttt{SimpleImputer(mode)}$\rightarrow$\texttt{OneHotEncoder}]\\
$\rightarrow$ \texttt{SelectKBest(mutual\_info, k)} $\rightarrow$ \texttt{LGBMClassifier(class\_weight='balanced', ...)}
\end{center}

La semilla está fijada en 42 en todos los componentes, incluidos la información mutua, que es estocástica, y el \textit{bootstrap}. La ejecución completa tarda unos 2 minutos en CPU. El notebook guarda el pipeline entrenado (\texttt{.pkl}), todas las métricas (\texttt{results.json}) y las figuras de este informe, por lo que cada cifra aquí reportada es reproducible.

%=====================================================================
\section{Resultados y análisis crítico}
\subsection{Desempeño en el conjunto de test (n = 1,518)}

\begin{figure}[H]
\centering
\begin{minipage}[c]{0.52\linewidth}
\centering\small
\captionof{table}{Métricas en test (umbral 0.5)}
\setlength{\tabcolsep}{5pt}
\begin{tabular}{lcc}
\toprule
Métrica & Valor & IC 95\,\% \\
\midrule
ROC-AUC & \textbf{0.979} & [0.973, 0.985] \\
PR-AUC (\textit{baseline} 0.362) & \textbf{0.960} & [0.945, 0.972] \\
Balanced accuracy & 0.929 & [0.915, 0.942] \\
Accuracy & 0.932 & [0.918, 0.944] \\
Precision (CONFIRMED) & 0.894 & [0.864, 0.920] \\
Recall (CONFIRMED) & 0.920 & [0.898, 0.942] \\
F1 (CONFIRMED) & 0.907 & [0.887, 0.924] \\
Especificidad & 0.938 & [0.922, 0.953] \\
\bottomrule
\end{tabular}
\end{minipage}\hfill
\begin{minipage}[c]{0.36\linewidth}
\centering
\includegraphics[width=\linewidth]{fig_confusion.png}
\captionof{figure}{Matriz de confusión en test.}
\end{minipage}
\end{figure}

\begin{figure}[H]
\centering
\includegraphics[width=0.8\linewidth]{fig_roc_pr.png}
\caption{Curvas ROC y Precision-Recall en test. La PR-AUC (0.960) está muy por encima de la proporción de planetas (0.362).}
\end{figure}

Las métricas son coherentes con la matriz: $P=506/566=0.894$ y $R=506/550=0.920$. El modelo pierde 44 de 550 planetas y acepta 60 de 968 falsos positivos. Como la prevalencia es del 36\,\%, la accuracy sola sería engañosa. Por eso se priorizan la PR-AUC, el recall de la clase CONFIRMED y la balanced accuracy. En test, XGBoost obtiene un AUC de 0.980 y LightGBM de 0.979. La diferencia pareada es $-0.0004$ \ci{-0.0019}{+0.0012}, y como el intervalo contiene el cero, ambos algoritmos son \textbf{estadísticamente indistinguibles}.

\subsection{Interpretabilidad}
En el modelo final, las variables más importantes (ganancia) son el radio planetario (0.20), la SNR del tránsito (0.13), el período (0.12), la duración (0.08), el parámetro de impacto (0.08) y la profundidad (0.06) (Fig.~\ref{fig:imp}, derecha). Este orden coincide con la física de los falsos positivos. Un radio inferido mayor que $\sim$15\,$R_\oplus$ delata una estrella compañera. Un parámetro de impacto alto con tránsito en forma de V es típico de las binarias rasantes. Una SNR baja aumenta la probabilidad de que la señal sea ruido o variabilidad estelar.

\begin{figure}[H]
\centering
\includegraphics[width=0.88\linewidth]{fig_importance_compare.png}
\caption{Importancia relativa (gain) de las 10 variables principales. Izquierda: experimento con columnas de incertidumbre (en rojo); derecha: modelo final.}\label{fig:imp}
\end{figure}

\subsection{Experimentos de \textit{leakage}}
Se reentrenó el modelo final, con los mismos hiperparámetros, agregando las columnas excluidas:
\begin{itemize}
\item \textbf{Con las 20 columnas \texttt{\_err} utilizables:} el AUC sube a 0.987 \ci{0.983}{0.991}. La diferencia pareada es $+0.008$ \ci{+0.005}{+0.011}: \textbf{real, pero pequeña}. Además, las barras de error pasan a dominar el ranking (\texttt{koi\_steff\_err1} es la primera variable, Fig.~\ref{fig:imp}, izquierda). Esa ganancia no proviene de la física, sino de información sobre el seguimiento ya realizado, que no existirá para un candidato nuevo. Renunciar a 0.008 de AUC a cambio de un modelo válido para su uso real es la decisión correcta.
\item \textbf{Con \texttt{koi\_score} y \texttt{koi\_fpflag\_*}:} el AUC sube a \textbf{1.000}, e incluso \texttt{koi\_score} por sí solo, sin modelo, alcanza 0.981. Esas columnas son la decisión del catálogo con otro nombre.
\end{itemize}

\subsection{Análisis de errores y sesgos}
\begin{itemize}
\item \textbf{Relación señal-ruido.} Es el factor dominante. El recall de planetas es 0.93--0.95 con SNR $>20$, baja a 0.89 con SNR entre 10 y 20, y a 0.44 con SNR $<10$ (aunque este último grupo solo tiene 9 planetas en test). Los falsos negativos tienen una mediana de SNR de 21, frente a 32 en los verdaderos positivos.
\item \textbf{Planetas pequeños.} Tanto los FN como los FP son objetos pequeños (mediana de 1.7\,$R_\oplus$) y de tránsito somero (300\,ppm), una región donde planetas reales e impostores se parecen físicamente. Los FP descartados correctamente tienen una mediana de 16\,$R_\oplus$: son binarias evidentes. Hay, por tanto, un \textbf{sesgo hacia planetas grandes y de alta SNR}: los análogos terrestres, los más interesantes científicamente, son los peor clasificados.
\item \textbf{Magnitud estelar.} Su efecto es despreciable: el recall va de 0.91 a 0.93 en todos los rangos de $K_p$, así que la SNR, y no el brillo de la estrella, explica los errores.
\item \textbf{Sesgo de selección en las etiquetas.} Los KOIs confirmados son los que fue factible observar. El modelo hereda ese sesgo, y su desempeño sobre los \texttt{CANDIDATE} restantes, más difíciles por construcción, probablemente sea menor que el medido en test.
\item \textbf{Sobreajuste.} Las curvas de aprendizaje muestran un AUC de entrenamiento de 1.000 frente a 0.982 en validación (brecha de 0.018). La curva de validación aún crece con más datos, lo que indica varianza moderada y que más ejemplos etiquetados seguirían ayudando.
\end{itemize}

\subsection{Uso operativo}
Aplicado al test, el umbral elegido \textit{out-of-fold} ($P\geq0.309$) recupera el \textbf{94.0\,\%} \ci{91.9}{95.9} de los planetas, con una precisión de 0.867 \ci{0.838}{0.893}, y solo señala el \textbf{39.3\,\%} \ci{36.9}{41.8} de los KOIs para seguimiento. Eso supone un \textbf{ahorro aproximado del 61\,\%} de tiempo de telescopio frente a observar todos. El recall en test (94\,\%) queda ligeramente por debajo del objetivo del 95\,\%, como era de esperar al aplicar a datos nuevos un umbral fijado de antemano, pero el objetivo cae dentro del intervalo de confianza. Aplicado a los 1,977 \texttt{CANDIDATE}, el modelo señala 747 para seguimiento, de los cuales 200 tienen $P>0.98$.

%=====================================================================
\section{Limitaciones y propuestas de mejora}
\begin{enumerate}
\item \textbf{Leakage residual posible:} \texttt{koi\_prad}, \texttt{koi\_srad} y \texttt{koi\_slogg} se derivan de parámetros estelares que también pueden refinarse tras la confirmación. Para descartarlo, habría que entrenar con los valores de una versión congelada del catálogo, por ejemplo el DR25 \cite{thompson2018}, anterior a la confirmación de cada objeto.
\item \textbf{Calibración de probabilidades} (\texttt{CalibratedClassifierCV}, isotónica) para que $P(\text{CONFIRMED})$ pueda usarse directamente en decisiones de costo-beneficio de tiempo de telescopio.
\item \textbf{Mejorar el desempeño en baja SNR y radios pequeños:} ponderar ejemplos por SNR, entrenar con las curvas de luz crudas mediante CNN 1D (que capturan la forma del tránsito, el ingreso/egreso y los eclipses secundarios) y explicar cada predicción con SHAP.
\item \textbf{Generalización:} validar en K2 y TESS, que tienen sistemáticos instrumentales distintos, con \textit{domain adaptation} o \textit{fine-tuning}. Usar \textit{active learning} para incorporar progresivamente los candidatos que se vayan resolviendo.
\end{enumerate}

%=====================================================================
\section{Conclusiones}
\begin{enumerate}
\item Se formuló la priorización de KOIs de Kepler como una clasificación binaria supervisada y se implementó un pipeline reproducible en el que todas las decisiones se tomaron sin usar el test. El modelo final alcanza \textbf{ROC-AUC 0.979} \ci{0.973}{0.985}, \textbf{PR-AUC 0.960} y \textbf{F1 0.907} en test.
\item El \textit{gradient boosting} superó claramente a la regresión logística (+0.05 de AUC), lo que confirma la naturaleza no lineal del problema. Entre LightGBM y XGBoost no hay diferencia significativa.
\item El control de leakage fue decisivo. Las variables de vetting llevan el AUC a 1.0 de forma espuria, y las barras de error aportan +0.008 de AUC que solo reflejan el seguimiento ya realizado. Sin ellas, el modelo se apoya en variables físicas: radio, SNR, período, duración y parámetro de impacto.
\item El principal punto débil son los planetas pequeños de baja SNR, que son precisamente los de mayor interés científico. Se propusieron mejoras concretas para ese caso.
\item En la práctica, el clasificador permite observar solo el $\sim$39\,\% de los KOIs (un ahorro de $\sim$61\,\%) recuperando el 94\,\% de los planetas, y entrega una lista priorizada de candidatos.
\end{enumerate}

\noindent\textbf{Código y notebook:} \url{https://github.com/Cha0smagick/Trabajo-unidad-IVmaestria-inteligencia-artificial-exoplanetas}

\vspace{-4pt}
{\small
\begin{thebibliography}{9}\setlength{\itemsep}{0pt}
\bibitem{borucki2010} Borucki, W.\,J. et al. (2010). Kepler Planet-Detection Mission: Introduction and First Results. \textit{Science}, 327(5968), 977--980.
\bibitem{thompson2018} Thompson, S.\,E. et al. (2018). Planetary Candidates Observed by Kepler. VIII. A Fully Automated Catalog. \textit{ApJS}, 235(2), 38.
\bibitem{xgboost} Chen, T. \& Guestrin, C. (2016). XGBoost: A Scalable Tree Boosting System. \textit{Proc. KDD}, 785--794.
\bibitem{lightgbm} Ke, G. et al. (2017). LightGBM: A Highly Efficient Gradient Boosting Decision Tree. \textit{NeurIPS}, 3146--3154.
\bibitem{sklearn} Pedregosa, F. et al. (2011). Scikit-learn: Machine Learning in Python. \textit{JMLR}, 12, 2825--2830.
\bibitem{nasa} NASA Exoplanet Archive. Kepler Objects of Interest — Cumulative Table. \url{https://exoplanetarchive.ipac.caltech.edu/} (consultado el 7 de noviembre de 2026).
\end{thebibliography}}

\end{document}
